%!TEX root = ../dissertation.tex

\chapter{Conclusion}
\label{chapter:conclusion}

This thesis originates from the key issue in synthetic data reconstruction that reconstructed data is only useful if its reliability can be proven. MOSA integrates multi-omic cancer cell line data and reconstructs missing omic profiles with integrated explainability, but accuracy and explainability constitute only part of what is needed for reliable use of this data. A clear and calibrated indication of confidence complements the model by ensuring the reliability of its reconstructions. In a biomedical setting, where reconstructed measurements may be used to guide hypotheses or prioritise further experiments, the absence of uncertainty information limits how safely such outputs can be interpreted. The main objective of this work was therefore to evaluate whether uncertainty estimation methods could provide useful information about the reliability of MOSA's reconstructions across different omic layers and reconstruction settings.

The results show that uncertainty estimates are informative but not uniformly reliable. Latent Space Sampling and Monte Carlo Dropout capture different aspects of predictive uncertainty and provide distinct conclusions. Latent Space Sampling, used here to analyse aleatoric uncertainty, generally produced narrow and homogeneous predictive distributions. The variability encoded in the VAE's latent representation is limited and aleatoric uncertainty of the samples propagated through the decoder is much smaller than the observed reconstruction errors. Out-of-sample, these intervals widened four- to five-fold for all continuous omics except CRISPR-Cas9, whose intervals widened only slightly, but they remained far narrower than the reconstruction errors. Monte Carlo Dropout, used as an approximation of epistemic uncertainty, produced substantially wider intervals in-sample for every continuous omic, although their magnitude depends on the chosen dropout configuration. As the two studied methods serve only as proxies for the two types of uncertainty, each not completely separating the other uncertainty effects, and as both methods are affected by external factors (i.e. modelling choices, encoder denoising and dropout rate selection), it can be concluded that model-related variability under Monte Carlo Dropout is larger in-sample than decoded latent space variability. Further study will be needed to evaluate whether aleatoric variability is intrinsically small or if the model's encoder fails to capture the data's heterogeneity. Similarly, further experiments are needed to determine whether the larger epistemic values are due to an aggressive regularisation setting.


A second important finding is that larger uncertainty estimates did not always imply better calibrated predictions. All continuous omics showed signs of overconfidence out-of-sample. Uncertainty should increase when less information is available for reconstruction. Under Latent Space Sampling it did, with intervals four to five times wider out-of-sample for all continuous omics except CRISPR-Cas9, whose intervals widened only slightly, but they remained far narrower than the reconstruction errors. Under Monte Carlo Dropout the opposite was observed for every continuous omic, most strongly for Methylation and CRISPR-Cas9, whose intervals narrowed by about 95\% or more. The Copy number reconstructions had a low in-sample calibration error under Monte Carlo Dropout, with a mean ECE of 0.077, but their calibration deteriorated out-of-sample, to an ECE of 0.264, and their mean Log Score rose from 1.31 to 5.15. These differences suggest that uncertainty in multi-omic reconstruction depends on omic-specific factors such as missingness and reconstruction difficulty. In Proteomics, features with more missing values were consistently more miscalibrated, and out-of-sample the features that were reconstructed worst tended to be the most miscalibrated.

This work's contribution is a structured evaluation of two accessible uncertainty estimation strategies in the specific context of VAE-based multi-omic reconstruction. By comparing aleatoric and epistemic uncertainty across omics and reconstruction settings, the analysis clarifies where uncertainty estimates can support interpretation and where they should be treated more cautiously.

This thesis also contains several limitations that should be taken into account. The analysis is restricted to the MOSA model and to the datasets considered in \cite{Cai2024}. This means that the conclusions may not generalise to other architectures or datasets. The uncertainty methods considered depend on implementation choices such as the number of samples and the dropout configuration. In particular, the dropout percentage used for Monte Carlo Dropout affects the magnitude of the output variance, while ECE and ENCE depend on the choice of calibration bin numbers and strategy. The difference in the number of observations used for ECE and ENCE may also affect the resulting estimates, since the bins can contain different numbers of samples between in-sample and out-of-sample reconstructions. The parameter values used for Monte Carlo Dropout, ECE and ENCE were selected based on common values in the literature. The evaluation is based on reconstruction performance and statistical uncertainty metrics, rather than on downstream biological validation, so it cannot determine whether uncertain reconstructions lead to incorrect biological conclusions. Lastly, the distinction between aleatoric and epistemic uncertainty is operational rather than exact. Latent Space Sampling and Monte Carlo Dropout provide useful proxies, but they do not fully separate the sources of uncertainty present in the data and model.

\section{Future Work}

This thesis focused on the analysis of the uncertainty of a single deep learning model. From this work, the logical next step would be to extend this evaluation beyond MOSA by evaluating other state of the art biological data reconstruction models. Useful baselines include MOFA+ for probabilistic multi-modal factor analysis \cite{Argelaguet2020MOFAplus}, mixOmics as a classic multivariate integration framework \cite{Rohart2017mixOmics}, and MOVE as a variational-autoencoder approach for multi-omics integration \cite{Allesoe2023MOVE}.
This benchmark would help better understand whether the patterns observed in this thesis are specific to MOSA or are properties of multi-omic reconstruction methods.

The conclusions of this work also suggest a route for improving MOSA. As seen in this thesis, the model's uncertainty estimates are often overconfident and are thus unreliable as measures of reconstruction quality. To improve this, future work should also focus on methods to improve model calibration or provide more robust uncertainty estimates. One first method for improving model calibration is post-hoc global uncertainty scaling, where predicted variances are adjusted on a validation set so that expected errors match the predicted uncertainty \cite{Levi2022RegressionCalibration}. Another method would be non-parametric post-hoc recalibration using quantile calibration which may help correct predictive distributions without the use of a single global scaling factor \cite{Kuleshov2018}. To provide more robust epistemic uncertainty estimates and reduce overconfidence in out-of-sample reconstructions, an ensemble method could additionally be explored by combining MOSA variants with different architectures, dropout configurations or asymmetric encoder-decoder designs.